\songtitle{From the Fields to the Flood}

\tag*{2026}C\tag{roads}
\begin{notes}
Part of the geography cycles series, lyrics by SUNO, on summaries by Grok, from Claud collaborations with Carlos Thompson.
\end{notes}
\begin{style}
Trap with oriental-scale melodies, high-pitched female vocals, half-time bounce, sliding sub bass, crisp hats, and a dark sparse verse-to-anthem arc; verses stay tense and detail-rich, pre-chorus strips to voice and a plucked motif, chorus opens with stacked doubles and short chantable lines, bridge drops to filtered drums and a lone ad-libbed top line, with risers, reverse swells, and porcelain chime accents; bright-on-top, heavy-in-the-low-end mix, world, deep, rich, distinct, female vocals
\end{style}
\begin{lyrics}
\songpart{Verse 1}
Yellow River stains the map\\
North plain built from silt and math\\
Shang bones, old soil, deep roots\\
Fields feed kings in temple boots

Beijing on the edge of stone\\
Steppe wind testing every throne\\
Wall line cut through pass and ridge\\
Hold that gate, hold that bridge

Xi’an by the Wei, old spark\\
Silk Road dust in the market arc\\
Persian thread in the woven cloth\\
Indian gold and a monk’s oath

\songpart{Pre-Chorus}
Three rivers, one world\\
Folded in the mud\\
Old roads, new trade\\
From the fields to the flood

\songpart{Chorus}
River circuit, run it through\\
Yellow, Yangtze, Seoul too\\
River circuit, hold that line\\
Land and sea in the same design

River circuit, feel the pull\\
Marsh and mountain, heavy and full\\
River circuit, side by side\\
One long pulse, one moving tide

\songpart{Verse 2}
Yangtze wide like a silver blade\\
Farm belts feed what the empire made\\
Shanghai where the harbor mouths\\
Pull the inland riches south

Delta rice and merchant hands\\
Tea ship sails to foreign lands\\
South China Sea, Pearl River glow\\
Hong Kong deep where the trade winds go

Seas keep turning the outside in\\
Salt on the edge, cash in the spin\\
Monsoon roads with the painted bowls\\
Ceramic dreams on African coasts

\songpart{Pre-Chorus}
Three rivers, one world\\
Folded in the mud\\
Old roads, new trade\\
From the fields to the flood

\songpart{Chorus}
River circuit, run it through\\
Yellow, Yangtze, Seoul too\\
River circuit, hold that line\\
Land and sea in the same design

River circuit, feel the pull\\
Marsh and mountain, heavy and full\\
River circuit, side by side\\
One long pulse, one moving tide

\songpart{Bridge}
Korea sits with the sea on three sides\\
Bridge and shield where the history hides\\
Han River basin, Seoul on the rise\\
Narrow waist with the watchful eyes

Writing, law, and the old belief\\
Passed to Japan like a burning leaf\\
Then the DMZ like a frozen scar\\
Thirty-eight line, visible from far

Taiwan on the strait, close but apart\\
Same deep current, different heart\\
China, Japan, Korea in view\\
One sea carries the old and new

\songpart{Chorus}
River circuit, run it through\\
Yellow, Yangtze, Seoul too\\
River circuit, hold that line\\
Land and sea in the same design

River circuit, feel the pull\\
Marsh and mountain, heavy and full\\
River circuit, side by side\\
One long pulse, one moving tide

\songpart{Outro}
One plain, one delta, one shore\\
Three zones, but the sea is more\\
River circuit, ring of the age\\
Written in mud on every page
\end{lyrics}

